\chapter{Implementation}\label{chapter:implementation}

\section{Pinning the dense weights}\label{sec:pinning}

The non-expert weights of the model file stay in the kernel-managed mmap region, and we mlock that region so the kernel cannot evict any of them.
The mlock excludes the experts (managed by the buffer of \autoref{sec:expert-buffer}) and the input embedding (each generated token reads only one of its rows).
Attention, the output projection, the layer norms, and the per-layer routers are therefore guaranteed resident throughout decode.


\section{The expert buffer}\label{sec:expert-buffer}

We replace the kernel's page-fault path for expert weights with \texttt{io\_uring}~\cite{axboe2019iouring} reads against an \texttt{O\_DIRECT}~\cite{open_man} file descriptor (\autoref{sec:bg-iouring}).
The application writes read requests into a shared submission ring, the kernel performs them in the background, and \texttt{O\_DIRECT} keeps the data out of the page cache so it transfers directly from the device into application buffers.
Together they let us submit many disk reads at once and pick up completions when convenient, with no intermediate copy in the page cache.

\texttt{O\_DIRECT} requires every read offset, length, and destination buffer address to be a multiple of 512 bytes.
The expert weight tensors in a GGUF file do not satisfy this.
We address the alignment with a one-time tool that reads every expert weight slice out of the GGUF and writes it into a sidecar file (\texttt{.bscexp}) at a 512-byte-aligned offset.
The sidecar is opened once per process with \texttt{O\_DIRECT}, and every \texttt{io\_uring} read targets a sidecar offset that is aligned by construction.

The buffer is a single 512-aligned anonymous region holding $N$ slots of 4.20~MiB each, where $N$ is configurable.
A hash map keyed by $(\mathit{layer}, \mathit{projection}, \mathit{expert\_id})$ returns the slot index of any resident slice.
On a hit, the buffer returns the slot's address.
On a miss, it picks a victim slot under the active replacement policy, issues a single \texttt{io\_uring} read into that slot, and returns the slot's address once the completion arrives.

Three replacement policies are implemented: LRU, LFU, and LFU-aging (the policies the simulator selected in \autoref{sec:tracer-sim}).
LRU keeps a doubly-linked list and promotes the slot on every access.
On a miss, the list's tail is evicted.
LFU and LFU-aging keep per-slot frequency counters and evict the slot with the smallest counter, with LFU-aging additionally halving every counter once per $3 \times \mathit{cache\_slots}$ accesses to discount stale popularity.
A per-load epoch prevents two misses in the same load from picking the same victim slot.


\section{Connecting the buffer to compute}\label{sec:buffer-integration}

The buffer's slices live in our own region of memory, separate from the mmap'd region of the model file that GGML's matrix-multiplication code normally reads from.
To make the slices visible to compute without copying, we use an existing GGML field: every tensor carries an opaque \texttt{extra} pointer that any computation in GGML can read freely.
Before each MoE layer runs, we fill \texttt{extra} with our array of slot pointers (one per active expert), and we patch the MoE matrix-multiplication function so that when \texttt{extra} is set, each active expert's address is taken from the array instead of from the tensor's base.
When \texttt{extra} is unset, the original code runs and no other operation in the model is affected.


\section{Pipelining I/O with compute}\label{sec:pipelining}

The simplest mode of the buffer is fully synchronous: a pre-compute callback before each layer submits all 12 reads as one \texttt{io\_uring} batch and waits for every completion before compute begins.
The three modes below progressively overlap I/O with compute, and each produces the same output as the synchronous path at fixed seed.


\subsection{Projection-group overlap}\label{sec:pipelining-po}

The down matmul takes the SwiGLU output as input, which means the down weights are not needed until SwiGLU has finished.
They can be loaded in the background while up, gate, and SwiGLU run.

We split the buffer's load into two phases inside a single callback before the up matmul.
Phase 1 submits the 8 up+gate reads as one \texttt{io\_uring} batch and waits for every completion.
Phase 2 submits the 4 down reads as a separate batch and returns immediately.
While the worker threads run up, gate, and SwiGLU on the up+gate slices, the 4 down reads complete in the background.
A second callback attached to the down matmul waits for any remaining down completions before down runs.

The 4 experts are still processed in a fixed order: each MoE matmul iterates them as $0, 1, 2, 3$ in turn (the positions in which the router emitted them).
Only the I/O is split.
The two phases share a single load epoch, which guarantees that phase 1's slots cannot be picked as eviction victims by phase 2.


\subsection{The fused MoE op}\label{sec:pipelining-fused}

In default llama.cpp, one MoE layer's feed-forward is a chain of separate GGML nodes: matmuls for up, gate, and down, plus SwiGLU and a per-expert accumulation.
Each matmul iterates the 4 selected experts internally.
The worker pool synchronizes at every node boundary, so all 4 experts finish their up matmul before any starts gate, and so on.
The structure is per-operation outside, per-expert inside.

Projection-group overlap fits this layout, because its callbacks attach to the existing nodes.
What the layout does not allow is starting one expert's down matmul while another expert is still loading from disk.
That would require inverting the loop nesting to per-expert outside, per-operation inside, and the per-node barriers prevent it.

We replace all of those default nodes with one fused custom GGML op.
From GGML's perspective, the entire MoE feed-forward is now one node, so GGML imposes no barrier inside it.
Inside, thread 0 picks the next compute action and issues any \texttt{io\_uring} submissions or waits the action requires.
All threads cross a barrier and cooperate on the picked action.
One barrier per dispatch round is enough, and the granularity at which we interleave I/O and compute is now ours to choose.

\begin{figure}[htbp]
\centering
\resizebox{0.95\textwidth}{!}{% Figure: GGML graph transformation for the fused MoE op.
% Compares the default graph (5 separate nodes) with our fused op (one node).
%
% Default (left): each MoE FFN layer is a chain of 5 GGML nodes. Each matmul
% node internally iterates over the 4 selected experts, and the worker pool
% synchronizes at every node boundary, so every expert finishes the current
% operation before any starts the next.
%
% Fused (right): the same compute is collapsed into one custom GGML node.
% Inside the node, our own dispatch loop runs the io_uring submissions, picks
% the next compute action under a state machine, and synchronizes the worker
% pool with an atomic barrier of our own.

\begin{tikzpicture}[
  font=\footnotesize,
  default_node/.style={
    draw=TUMSecondaryBlue!70!black, fill=TUMSecondaryBlue!12,
    rounded corners=2pt,
    minimum width=28mm, minimum height=7mm,
    align=center
  },
  fused_outer/.style={
    draw=TUMAccentGreen!75!black, fill=TUMAccentGreen!5,
    line width=0.8pt, rounded corners=4pt,
    minimum width=44mm, minimum height=46mm
  },
  fused_inner/.style={
    draw=TUMAccentGreen!60!black, fill=TUMAccentGreen!18,
    rounded corners=2pt,
    minimum width=36mm, minimum height=7mm,
    align=center, font=\scriptsize
  },
  arr/.style={->, line width=0.45pt, draw=black!75},
  hdr/.style={font=\sffamily\bfseries\small, align=center},
  foot/.style={font=\scriptsize, text=black!65, align=center},
]

  % --- LEFT column: default graph ---
  \node[hdr] at (0, 7.0) {Default GGML graph};

  \node[default_node] (up_d)     at (0, 5.5) {matmul (up)};
  \node[default_node] (gate_d)   at (0, 4.4) {matmul (gate)};
  \node[default_node] (swiglu_d) at (0, 3.3) {SwiGLU};
  \node[default_node] (down_d)   at (0, 2.2) {matmul (down)};
  \node[default_node] (acc_d)    at (0, 1.1) {accumulate};

  \draw[arr] (up_d)     -- (gate_d);
  \draw[arr] (gate_d)   -- (swiglu_d);
  \draw[arr] (swiglu_d) -- (down_d);
  \draw[arr] (down_d)   -- (acc_d);

  \node[foot] at (0, 0.0)
    {each matmul iterates 4 experts internally.\\Worker pool synchronizes at every node boundary.};

  % --- middle: replacement arrow ---
  \node[font=\Large] at (5.6, 3.3) {$\Rightarrow$};

  % --- RIGHT column: fused op ---
  \node[hdr] at (9.0, 7.0) {Fused custom GGML op};

  \node[fused_outer] (fused) at (9.0, 3.3) {};

  \node[fused_inner] (sm)      at (9.0, 5.0) {state-machine picker};
  \node[fused_inner] (iouring) at (9.0, 4.0) {\texttt{io\_uring} submit / wait};
  \node[fused_inner] (compute) at (9.0, 3.0) {up, gate, SwiGLU, down};
  \node[fused_inner] (acc_f)   at (9.0, 2.0) {accumulate};

  \draw[arr] (sm)      -- (iouring);
  \draw[arr] (iouring) -- (compute);
  \draw[arr] (compute) -- (acc_f);

  \node[foot] at (9.0, 0.0)
    {one node from GGML's view.\\Internal dispatch under our control.};

\end{tikzpicture}
}
\caption[Default GGML graph versus the fused MoE op]{Default GGML graph (left) versus the fused custom op (right). The 5 separate nodes on the left are subject to GGML's per-node barriers, which fix the loop nesting to per-operation outside, per-expert inside. Collapsing them into one node on the right places the dispatch loop inside our control, so the per-expert and per-matmul dispatch in the next two subsections become possible.}
\label{fig:moe-fused-op}
\end{figure}


\subsection{Async-projection-overlap}\label{sec:pipelining-apo}

Async-projection-overlap differs from projection-group overlap in two ways.
All 12 reads are submitted at the start of the fused op, not split into two phases.
The 4 experts are processed in the order their data arrives, not in $0, 1, 2, 3$ position order.

Each expert's up and gate reads share one \texttt{io\_uring} tag and its down read takes a separate tag, for $4 \times 2 = 8$ tags in flight together.
Thread 0 waits for any one of the 4 up+gate pairs to complete and dispatches that expert.
The worker threads run up, gate, and SwiGLU on that expert.
Thread 0 then waits for that expert's down read (in flight since layer entry) before the worker threads run down.
Thread 0 picks the next first-ready expert from the remaining tags and the loop repeats until all 4 have run.

Each expert's output is written into a slot at its router-assigned position, and the 4 slots are summed in position order at the end of the op, so the final output matches the synchronous and projection-group overlap paths.


\subsection{Async-experts}\label{sec:pipelining-ae}

Async-experts pushes the dispatch granularity one step finer.
Where async-projection-overlap dispatches per expert (each expert runs up + gate + SwiGLU + down end-to-end before the next expert starts), async-experts dispatches per individual matmul: an expert's up matmul can run as soon as its up tag completes, even before its gate read has arrived.

Each projection of each expert gets its own \texttt{io\_uring} tag, for $4 \times 3 = 12$ tags in flight together.
Thread 0 maintains a per-expert state machine and at each dispatch round walks the 4 active experts in order, picking the first eligible action (\autoref{fig:async-experts-state}).
When the second of an expert's two pre-SwiGLU matmuls completes, SwiGLU and the down-input quantization fold into the same dispatch round, so SwiGLU is never a standalone action.

\begin{figure}[htbp]
\centering
\resizebox{0.85\textwidth}{!}{% Figure: per-expert state machine for the async-experts dispatcher.
% Matches the actual implementation in llama-moe-pipeline.cpp:
%   AE_UP_MATMUL        : up only (gate not yet computed)
%   AE_UP_AND_FINISH    : up + swiglu + quantize (gate already computed)
%   AE_GATE_MATMUL      : gate only (up not yet computed)
%   AE_GATE_AND_FINISH  : gate + swiglu + quantize (up already computed)
%   AE_DOWN_MATMUL      : down only
%
% The fold-in optimization (AND_FINISH variants) collapses SwiGLU and the
% down-input quantize into the second-arrival matmul, so there is no
% standalone "both matmuls done, no swiglu" intermediate state.
%
% Conventions:
%   * S_0 (green)  : initial, nothing computed
%   * S_1          : up matmul done first
%   * S_2          : gate matmul done first
%   * S_3          : SwiGLU + quantize done; ready for down
%   * S_4 (orange) : down matmul done; expert complete
%   * Edge labels  : action above (blue), io_uring completion guard below
%                    (italic orange)

\begin{tikzpicture}[
    font=\sffamily\small,
    >={Latex[length=2.5mm]},
    state/.style={
        draw, rounded corners=2pt, fill=TUMLightGray!20,
        minimum width=18mm, minimum height=11mm, align=center,
        inner sep=2pt,
        font=\sffamily\small,
    },
    initstate/.style={state, fill=TUMAccentGreen!20, draw=TUMAccentGreen!60!black},
    finalstate/.style={state, fill=TUMAccentOrange!25, draw=TUMAccentOrange!70!black},
    actlabel/.style={font=\scriptsize\sffamily, text=TUMSecondaryBlue, inner sep=1.5pt},
    guardlabel/.style={font=\scriptsize\sffamily\itshape, text=TUMAccentOrange!85!black, inner sep=1.5pt},
]
  % --- States ---
  \node[initstate]  (s0) at (0,    0)    {$S_0$\\[1pt]\scriptsize init};
  \node[state]      (s1) at (4.5,  1.6)  {$S_1$\\[1pt]\scriptsize up done};
  \node[state]      (s2) at (4.5, -1.6)  {$S_2$\\[1pt]\scriptsize gate done};
  \node[state]      (s3) at (9.0,  0)    {$S_3$\\[1pt]\scriptsize SwiGLU done};
  \node[finalstate] (s4) at (13.5, 0)    {$S_4$\\[1pt]\scriptsize complete};

  % --- S_0 to first-arrival branches (diamond top + bottom) ---
  \draw[->, thick] (s0) -- (s1)
      node[actlabel,   sloped, above, pos=0.5] {up matmul}
      node[guardlabel, sloped, below, pos=0.5] {loaded\_up};
  \draw[->, thick] (s0) -- (s2)
      node[actlabel,   sloped, above, pos=0.5] {gate matmul}
      node[guardlabel, sloped, below, pos=0.5] {loaded\_gate};

  % --- Second-arrival fold-in: AND_FINISH variants merge to S_3 ---
  \draw[->, thick] (s1) -- (s3)
      node[actlabel,   sloped, above, pos=0.5, align=center] {gate matmul\\+ SwiGLU + quant}
      node[guardlabel, sloped, below, pos=0.5] {loaded\_gate};
  \draw[->, thick] (s2) -- (s3)
      node[actlabel,   sloped, above, pos=0.5, align=center] {up matmul\\+ SwiGLU + quant}
      node[guardlabel, sloped, below, pos=0.5] {loaded\_up};

  % --- Final transition to S_4 ---
  \draw[->, thick] (s3) -- (s4)
      node[actlabel,   above, pos=0.5] {down matmul}
      node[guardlabel, below, pos=0.5] {loaded\_down};
\end{tikzpicture}
}
\caption[Per-expert state machine for async-experts]{Per-expert state machine inside the async-experts dispatch loop. $S_0$ is the initial state. The diamond $S_0 \to \{S_1, S_2\} \to S_3$ captures the difference from async-projection-overlap: an expert's first matmul runs as soon as one of its up or gate reads completes, without waiting for the other. The second-arrival matmul folds SwiGLU and the down-input quantization into the same dispatch round (the \emph{and finish} variants), so $S_3$ is the post-SwiGLU state ready for the down matmul.}
\label{fig:async-experts-state}
\end{figure}

The same position-order summation at the end makes the output identical to async-projection-overlap.
In our wall-clock measurements (\autoref{chapter:evaluation}), async-experts lands within roughly one percent of async-projection-overlap, so we keep async-projection-overlap as the deployed configuration.
